\chapter{Literature Review}
\label{ch:literature}

\section{Session-Based Recommendation}

Early session-based recommenders relied on item-to-item similarity: the item-kNN
approach recommends the items that most often co-occur in sessions with the
current item. It remains a strong baseline.

\textbf{Hidasi et al.\ (2016)} \cite{hidasi2016} introduced GRU4Rec, the first
RNN for session-based recommendation. Its central engineering contribution is
\emph{session-parallel mini-batches}: each mini-batch slot holds one session and
advances one click per step; when a session ends, a new session takes its slot
and the slot's hidden state is reset. On YooChoose, GRU4Rec clearly outperformed
item-kNN.

\textbf{Hidasi and Karatzoglou (2018)} \cite{hidasi2018} improved GRU4Rec with
additional negative samples shared across the mini-batch and drawn in proportion
to popularity$^\alpha$, and with the BPR-max and TOP1-max ranking losses, raising
Recall@20 on the full YooChoose data from 0.585 to 0.721. This improved version
is the one used in this project.

\textbf{Hidasi and Czapp (2023)} \cite{hidasi2023} examined six third-party
implementations of GRU4Rec used in comparative studies and found that they
scored 30--99\% below the official implementation, because of incorrect loss
functions, wrong hidden-state handling, missing negative sampling, or broken
dropout. Many papers that report GRU4Rec as a weak baseline used such
implementations. For this reason all session-based models in this project are
built on the official PyTorch code.

\textbf{Li et al.\ (2017)} \cite{li2017} proposed NARM, which adds an attention
mechanism over the GRU hidden states to capture the user's main purpose in the
session. \textbf{Liu et al.\ (2018)} \cite{liu2018} proposed STAMP, an
attention-based model of short-term interest. \textbf{Wu et al.\ (2019)}
\cite{wu2019} modelled sessions as graphs with SR-GNN. These papers established
the standard yoochoose1/64 and yoochoose1/4 benchmarks used in
Chapter~\ref{ch:results}.

\textbf{Kang and McAuley (2018)} \cite{kang2018} proposed SASRec, a
self-attentive sequential recommender, and \textbf{Sun et al.\ (2019)}
\cite{sun2019} the bidirectional BERT4Rec. Transformer-based models capture
long-range dependencies but are larger and more expensive than GRU4Rec.

\section{Next-Basket Recommendation}

\textbf{Rendle et al.\ (2010)} \cite{rendle2010} introduced FPMC, which combines
matrix factorisation with a personalised Markov chain over consecutive baskets.
\textbf{Yu et al.\ (2016)} \cite{yu2016} proposed DREAM, the first RNN for
next-basket recommendation, which pools the item embeddings of each basket and
feeds the pooled vectors to an RNN.

\textbf{Guidotti et al.\ (2019)} \cite{guidotti2019} proposed TARS, which mines
temporal annotated recurring sequences from each customer's own history and
predicts the next basket from them; it is interpretable by construction. Its
implementation, \texttt{tbp-next-basket}, also provides an evaluation suite of
set-based precision, recall, F-scores and hit rate.

\textbf{Hu and He (2019)} \cite{hu2019} proposed Sets2Sets, an encoder-decoder
with an explicit repeat component. \textbf{Hu et al.\ (2020)} \cite{hu2020}
showed with TIFU-KNN that a simple method based on time-decayed personal item
frequencies, combined with the frequencies of nearest-neighbour users,
outperforms most neural models. \textbf{Ariannezhad et al.\ (2022)}
\cite{ariannezhad2022} proposed ReCANet, a model dedicated to repeat
consumption.

\textbf{Li et al.\ (2023)} \cite{li2023} re-evaluated next-basket methods and
found that repeat items account for most of the correct predictions on Instacart
and Dunnhumby, that personal-frequency baselines are very strong, and that
several neural models perform close to global popularity. They recommend
reporting repeat and explore performance separately, which this project does.

\textbf{Van Maasakkers, Fok and Donkers (2023)} \cite{vanmaasakkers2023} applied
a GRU directly to multi-hot basket vectors, with a sigmoid output for every
product, on Instacart and Dunnhumby (The Complete Journey). With large hidden
layers (750 and 1500 units) the GRU outperformed personal frequency and TARS on
most metrics. Their preprocessing and evaluation protocol is followed here, and
their basket-GRU is the next-basket baseline that KAN layers are added to.

\section{Kolmogorov-Arnold Networks}

\textbf{Liu et al.\ (2025)} \cite{liu2025kan} introduced KANs, which place a
learnable univariate function, parameterised by B-splines, on every edge. On
small scientific and function-fitting tasks KANs reached a given accuracy with
fewer parameters than MLPs and allowed learned functions to be inspected and
even replaced by symbolic formulas. The authors identified slow training as the
main limitation.

Two later studies qualify these results. \textbf{Yu, Yu and Wang (2024)}
\cite{yu2024} compared KANs and MLPs with matched parameters and FLOPs across
machine learning, computer vision, NLP, audio and symbolic formula
representation; KANs were better only on symbolic formulas, and an MLP given
B-spline activations matched or exceeded KAN elsewhere.
\textbf{Hou et al.\ (2025)} \cite{hou2025} critically assessed the claims made
for KANs and reported that they are substantially slower to train than MLPs and
rarely more accurate outside symbolic regression. Both studies motivate the
parameter-matched MLP controls used in this project.

Several cheaper bases have been proposed. \textbf{FastKAN} \cite{li2024fastkan}
replaces B-splines with Gaussian radial basis functions and shows that a KAN
with such a basis is a radial basis function network. Chebyshev-polynomial KANs
\cite{ss2024cheby} use orthogonal polynomials on tanh-bounded inputs, and
ReLU-KAN \cite{relukan} builds the basis from ReLU products. The
\textbf{Kolmogorov-Arnold Transformer} \cite{yang2025kat} introduced group-rational
KAN (GR-KAN) layers, in which rational functions are shared within groups of
input channels, and used them as drop-in replacements for the MLP blocks of
vision transformers.

\section{KANs for Sequential and Temporal Data}

\textbf{Xu, Chen and Wang (2024)} \cite{xu2024tskan} applied KANs to time-series
forecasting and argued that the learned univariate functions make the models
interpretable. \textbf{Genet and Inzirillo (2024)} \cite{genet2024} proposed TKAN,
a recurrent cell that keeps the LSTM's sigmoid gates and memory cell and uses KAN
layers only inside the recurrent transformation. TKAN shows that KANs can be
placed inside a recurrent cell without giving up gating; the KAN-GRU cell of
Chapter~\ref{ch:models} follows the same principle.

\section{KANs in Recommender Systems}

\textbf{CF-KAN} \cite{park2024cfkan} replaces the MLP of an autoencoder-based
collaborative filtering model with a KAN and reports higher Recall@20 than strong
collaborative filtering baselines on MovieLens-1M, Yelp and Anime, together with
more robustness to catastrophic forgetting. It is, to our knowledge, the only KAN
recommender that compares against an MLP version of the same model with the same
number of parameters, and it is not yet peer-reviewed. \textbf{FourierKAN-GCF}
\cite{xu2024fkgcf} replaces the feature transformation in graph collaborative
filtering with a Fourier-basis KAN. \textbf{KANRec} \cite{xu2026kanrec} uses KAN
layers to fuse global and local interests in sequential recommendation.

\section{Summary of Related Work}


\begin{table}[H]
\centering
\caption{Summary of the main related work.}
\label{tab:related}
\small
\begin{tabular}{p{3.4cm} p{0.7cm} p{3.2cm} p{2.6cm} p{4.2cm}}
\toprule
\textbf{Work} & \textbf{Year} & \textbf{Model} & \textbf{Data} & \textbf{Relevance / limitation} \\
\midrule
Hidasi et al.\ \cite{hidasi2016,hidasi2018} & 2016, 2018 & GRU4Rec & YooChoose & Strong session baseline; reimplementations often flawed \cite{hidasi2023} \\
Li et al.\ \cite{li2017} & 2017 & NARM (GRU + attention) & YooChoose, Diginetica & Standard benchmark; black-box attention \\
Wu et al.\ \cite{wu2019} & 2019 & SR-GNN & YooChoose, Diginetica & Standard yoochoose1/64 table \\
Yu et al.\ \cite{yu2016} & 2016 & DREAM (RNN) & Ta-Feng, T-Mall & Close to popularity in later re-evaluations \cite{li2023} \\
Guidotti et al.\ \cite{guidotti2019} & 2019 & TARS & Coop, Ta-Feng & Interpretable; per-customer, slow \\
Hu et al.\ \cite{hu2020} & 2020 & TIFU-KNN & Instacart, Dunnhumby & Simple frequency model beats neural models \\
Li et al.\ \cite{li2023} & 2023 & Re-evaluation & Instacart, Dunnhumby, Ta-Feng & Repeat items dominate accuracy \\
van Maasakkers et al.\ \cite{vanmaasakkers2023} & 2023 & Basket-GRU & Instacart, Dunnhumby & Our basket baseline and protocol \\
Liu et al.\ \cite{liu2025kan} & 2025 & KAN & Function fitting, science & Slow training \\
Yu et al.\ \cite{yu2024} & 2024 & KAN vs MLP & Many domains & KAN rarely wins when matched \\
Genet and Inzirillo \cite{genet2024} & 2024 & TKAN & Time series & KAN inside a gated cell \\
Park et al.\ \cite{park2024cfkan} & 2024 & CF-KAN & ML-1M, Yelp, Anime & KAN in collaborative filtering \\
Xu et al.\ \cite{xu2026kanrec} & 2026 & KANRec & Sequential rec. & KAN interest fusion \\
\bottomrule
\end{tabular}
\end{table}

\section{Research Gap}

KAN layers have been applied to collaborative filtering, graph recommendation and
attention-based sequential recommendation, and KAN components have been placed
inside gated recurrent cells for time series. We found no work that places KAN
layers inside GRU4Rec, the most widely used RNN for session-based recommendation,
or that applies KANs to next-basket recommendation, and among the KAN
recommenders we found, only CF-KAN controls for model size with a
parameter-matched MLP. The first draft of this
project also showed how easily such comparisons go wrong: without a standard
protocol, a verified implementation and a capacity-matched control, neither a
gain nor a loss can be attributed to the KAN. This project therefore asks a
narrower question than ``are KANs better?'': \emph{where in a GRU-based
recommender, if anywhere, does a KAN layer help, compared with an MLP of the same
size, on standard session and basket benchmarks?}
